\section{Main results}

The minimax criterion in \cref{obj:def:criterion} admits a two-sided characterization in terms of the number of coordinate pairs relative to the sample size. For presentation, define the interaction count \(B\) and the comparison scale \leanref{sym:a_s}{\ensuremath{a_s(n,d)}} by
\[
B=\binom d2,
\qquad
a_s(n,d)=\min\{1,(B/n)^s\},
\qquad n,d\geq2,\quad 0<s\leq1.
\]
The following result establishes the log-free capacity bound and the exact dimension condition for vanishing excess. The reflection lift and Gaussian completion appearing in the statement are specified in \cref{obj:synth_2,obj:synth_4}.

\begin{definitionv}{synth_4}[Pre-assignment Gaussian completion]
The Gaussian completion is the conditional law of the pre-assignment potential-outcome array \(Y\) given \(X\), generated independently across units by
\[
Y_i(0)=m(X_i)-\frac{0.2+0.1\sqrt2\sin(2\pi X_{i1})}{2}+\varepsilon_{i0},
\qquad
Y_i(1)=m(X_i)+\frac{0.2+0.1\sqrt2\sin(2\pi X_{i1})}{2}+\varepsilon_{i1}.
\]
The variables \(\varepsilon_{ia}\) are mutually independent standard Gaussians, independent of the covariates, and \(m\) is the centered measurable square-integrable prognostic function supplied to the completion.
\end{definitionv}

The comparison scale separates two regimes. In the subcritical regime \(B<n\), minimax excess is of order \((B/n)^s\): the interaction count enters through its ratio to the number of experimental units. When \(B\geq n\), the scale saturates at one, and the displayed lower and upper bounds are constants independent of \(n\) and \(d\). The assignment procedure \leanref{sym:pi_star}{\ensuremath{\pi^\star_{n,d}}} specified in \cref{obj:thm:log-free-bivariate-capacity} attains this scale in both regimes, using independent fair assignment in the saturated regime and spectral rounding in the subcritical regime.

At the endpoint \(s=1\), the subcritical rate is \(B/n\), with the same upper constant and a strictly positive lower constant. More generally, the lower constant \leanref{sym:c_s}{\ensuremath{c_s}} is bounded below by \(c_1>0\) throughout \(0<s\leq1\). The two-sided comparison therefore describes the endpoint together with the fractional-smoothness cases, with a logarithm-free rate. For each fixed smoothness value, it also gives an exact criterion for vanishing minimax scaled excess along arbitrary integer dimension sequences: the number of coordinate pairs must be negligible relative to sample size, equivalently \(d_n=o(\sqrt n)\). This equivalence accommodates irregular as well as monotone dimension sequences.

\subsection{Simultaneous attainment}

For the spectral branch of the assignment procedure, the auxiliary randomization transforms the covariates of the original units as follows.

\begin{definitionv}{synth_2}[Randomized reflection lift]
Let \(\mathbb T_2^d=[0,2)^d\) with coordinatewise addition modulo two. For \(x\in[0,1]^d\) and \(\eta\in\{0,1\}^d\), define the coordinatewise reflection map \(\operatorname{lift}:[0,1]^d\times\{0,1\}^d\to\mathbb T_2^d\) by
\[
\operatorname{lift}(x,\eta)_j
=
\begin{cases}
x_j \pmod 2,&\eta_j=0,\\
2-x_j \pmod 2,&\eta_j=1.
\end{cases}
\]
For the original sample \(X=(X_1,\ldots,X_n)\), take independent fair coins \(\eta_{ij}\), independent of \(X\), the common shift, and all rounding seeds, and put
\[
\widetilde X_i=\operatorname{lift}(X_i,\eta_i),
\qquad
W_i=\widetilde X_i+T\pmod2,
\]
where \(\eta_i=(\eta_{i1},\ldots,\eta_{id})\) and \(T\) is an independent common Haar-uniform shift on \(\mathbb T_2^d\).
\end{definitionv}

The reflections act separately across units and coordinates, whereas the Haar shift is shared by all units. The transformed covariates in \cref{obj:synth_2} supply the spectral rows used to generate assignment signs for the original sample.

The simultaneous guarantee separates the choice of assignment from the smoothness value used to assess precision. The procedure in \cref{obj:thm:log-free-bivariate-capacity} uses the same spectral rows and randomization for every \(0<s\leq1\), while the resulting bound varies with the exponent \(s\). Its finite completion property applies to every original sample and every allowed realization of the auxiliary seeds. Thus the displayed risk bounds describe a fully specified assignment law on the original units.

\subsection{Log-free capacity and dimension frontier}

The lower bound has a separate statistical role: it applies to every admissible assignment law through a single prior chosen independently of the covariate sample. The prior specification and its class-membership properties are developed further in \cref{obj:def:cosine-prior,obj:lem:cosine-prior}.

The normalized prognostic function \leanref{sym:m_xi}{\ensuremath{m_\xi}} in \cref{obj:def:cosine-prior} ranges over a finite family of pair-cosine functions. Its law, \leanref{sym:nu_star}{\ensuremath{\nu^\star_{n,d,s}}}, depends on sample size, dimension, and smoothness, while remaining independent of both the covariate sample and the chosen assignment law. Every realization belongs to the centered pooled Sobolev class by \cref{obj:thm:log-free-bivariate-capacity}. The contraction inequality used in the supporting converse analysis is \citet[Theorem~12(4), with Rademacher complexity as in Definition~2]{BartlettMendelson2002Complexities}; its applications here use the absolute-value map with finite classes and finite-dimensional Euclidean unit-ball linear classes. The following result states the resulting prior-average bound over the full design class.

\begin{theoremv}{thm:log-free-bivariate-capacity}[Log-free capacity and dimension frontier]*
For every pair of integers \(n,d\geq2\), write
\[
B=\binom d2.
\]
There is a single fair, outcome-independent Borel original-sample assignment kernel \(\pi^\star_{n,d}\), independent of the smoothness parameter, belonging to the design class of \cref{obj:def:design-class}. It assigns independent fair signs when \(n\leq B\). When \(B<n\), it is the law of the signs produced by spectral rounding with independent reflection coins, an independent common Haar shift, and independent uniform rounding seeds.

More precisely, the reflected and shifted sample is
\[
W_i=\operatorname{lift}(X_i,\text{reflection coins})+T,
\]
where \(T\) is the common Haar shift on the period-two torus. The input rows are the first \(\lfloor n/4\rfloor\) prescribed real features evaluated at \(W_i\). For every choice of reflection coins and common shift, and every rounding-seed array whose entries belong to \([0,1]\), the fractional coordinate after \(n\) rounding iterations equals the corresponding output sign, for every original sample and every unit.

For every real \(s\) with \(0<s\leq1\), let
\[
a_s(n,d)=b_s(n,d)=\min\{1,(B/n)^s\},
\qquad
c_s=\frac{2}{(48+32\pi^2)16384^s},
\qquad
c_1=\frac{2}{(48+32\pi^2)16384}.
\]
Thus \(a_s(n,d)=1\) when \(n<B\), and \(a_s(n,d)=(B/n)^s\) when \(B\leq n\). The constants satisfy
\[
0<c_1\leq c_s<3072.
\]
With \(\mathcal H_{d,s}\) as in \cref{obj:def:sobolev-class}, and \(R_s(n,d)\) the minimax criterion of \cref{obj:def:criterion}, define the scaled signing loss by
\[
\mathcal L_{n,d,s}(\pi,m)
=\frac4n\,\mathbb E_{\pi}
 \left(\sum_{i=1}^n Z_i m(X_i)\right)^2.
\]
The same kernel satisfies, simultaneously for all \(0<s\leq1\),
\[
c_1a_s(n,d)
\leq c_sa_s(n,d)
\leq R_s(n,d)
\leq \sup_{m\in\mathcal H_{d,s}}
       \mathcal L_{n,d,s}(\pi^\star_{n,d},m)
\leq3072a_s(n,d).
\]

For each such \(s\), the finite coefficient prior \(\nu^\star_{n,d,s}\), with normalization
\[
C_0=\sqrt{48+32\pi^2},
\]
has every coefficient-sign realization \(m_\xi\) centered and in \(\mathcal H_{d,s}\):
\[
\int_{[0,1]^d}m_\xi(x)\,dx=0,
\qquad
m_\xi\in\mathcal H_{d,s}.
\]
Drawn independently of the covariate sample, it certifies, for every admissible design \(\pi\),
\[
c_sa_s(n,d)
\leq
\mathbb E_{m\sim\nu^\star_{n,d,s}}
       \mathcal L_{n,d,s}(\pi,m).
\]

Let \(\mathcal Q_{d,s}\) be the frozen uniform-covariate, zero-intercept, correctly specified outcome-law class with finite second moments and centered prognostic function in \(\mathcal H_{d,s}\). For a single-unit law \(P\), its prognostic function, actual HT estimator, law-specific benchmark, and scaled excess are
\[
m_P(u)=\mathbb E_P[(O(1)+O(0))/2\mid U=u],
\qquad
\widehat T_P=\frac2n\sum_{i=1}^n
 Z_iO_i((Z_i+1)/2),
\]
\[
\mathsf V(P)
=\operatorname{Var}_P\!\left(
 \mathbb E_P[O(1)-O(0)\mid U]\right)
+2\mathbb E_P\operatorname{Var}_P(O(1)\mid U)
+2\mathbb E_P\operatorname{Var}_P(O(0)\mid U),
\]
\[
\mathcal V_n(\pi,P)
=n\operatorname{Var}_{P,\pi}(\widehat T_P)-\mathsf V(P).
\]
Here the experiment uses independent original-unit copies, and conditional means are interpreted modulo covariate null sets. Then
\[
\inf_{\pi\in\mathcal D_{n,d,s}}
 \sup_{P\in\mathcal Q_{d,s}}\mathcal V_n(\pi,P)
=R_s(n,d),
\]
\[
\sup_{P\in\mathcal Q_{d,s}}
 \mathcal V_n(\pi^\star_{n,d},P)
=
\sup_{m\in\mathcal H_{d,s}}
 \mathcal L_{n,d,s}(\pi^\star_{n,d},m).
\]
The Gaussian completions of functions in \(\mathcal H_{d,s}\) form a proper subclass of \(\mathcal Q_{d,s}\). Restricting the law supremum to these Gaussian completions gives the same minimax capacity \(R_s(n,d)\) and the same worst-case loss for \(\pi^\star_{n,d}\).

For every pair of real density bounds satisfying
\[
0<\mathrm{lower}\leq1\leq\mathrm{upper},
\]
let \(\mathcal A_{d,s}\) be the bounded-density class defined in \cref{obj:lem:published-prognostic-capacity}: its laws have finite potential-outcome second moments, satisfy
\[
\int m_P(u)^2\,dP_U(u)\leq1,
\]
and have \(m_P\) equal almost everywhere to a permitted intercept plus a centered function in \(\mathcal H_{d,s}\), with covariate density between the displayed bounds. Its minimax excess satisfies
\[
R_s(n,d)
\leq
\inf_{\pi\in\mathcal D_{n,d,s}}
 \sup_{P\in\mathcal A_{d,s}}\mathcal V_n(\pi,P).
\]

For every admissible \(\pi\) and every \(m\in\mathcal H_{d,s}\), the actual HT estimator \(\widehat\theta\) under the Gaussian completion of \(m\), with benchmark \(V^*=4.01\), satisfies
\[
n\operatorname{Var}(\widehat\theta)-4.01
=\mathcal L_{n,d,s}(\pi,m)<\infty.
\]
Consequently, for every \(\varepsilon>0\),
\[
n\geq B\max\!\left\{
1,\left(\frac{3072}{4.01\varepsilon}\right)^{1/s}
\right\}
\quad\Longrightarrow\quad
\sup_{m\in\mathcal H_{d,s}}
 \mathcal L_{n,d,s}(\pi^\star_{n,d},m)
\leq4.01\varepsilon.
\]
Conversely, for every \(0<\varepsilon<c_s/4.01\),
\[
R_s(n,d)\leq4.01\varepsilon
\quad\Longrightarrow\quad
n\geq B\left(\frac{c_s}{4.01\varepsilon}\right)^{1/s}.
\]

Finally, for every fixed \(s\in(0,1]\) and every integer sequence \(d_n\geq2\),
\[
R_s(n,d_n)\longrightarrow0
\quad\Longleftrightarrow\quad
\frac{\binom{d_n}{2}}n\longrightarrow0
\quad\Longleftrightarrow\quad
d_n=o(\sqrt n).
\]
For every fixed pair \(0<\mathrm{lower}\leq1\leq\mathrm{upper}\), vanishing minimax excess over the corresponding bounded-density classes implies the same dimension frontier:
\[
\inf_{\pi\in\mathcal D_{n,d_n,s}}
 \sup_{P\in\mathcal A_{d_n,s}}\mathcal V_n(\pi,P)
\longrightarrow0
\quad\Longrightarrow\quad
d_n=o(\sqrt n).
\]
\end{theoremv}
% CAUSALSMITH-CITED-SCOPE-BEGIN thm:log-free-bivariate-capacity
\verificationfootnotetext{\textbf{Formalization scope.} This result uses the cited conclusion \citep{BartlettMendelson2002Complexities} (Theorem 12(4), with Rademacher complexity as in Definition 2) and \citep{Cytrynbaum2026Limits} (Proposition 2.3 (Excess Variance), equation (2.4), with Assumption 2.1 and Definition 2.2). The cited source proof is not formalized here: Lean verifies its use and all remaining steps. If that cited conclusion is false, the portions of this result that depend on it are not certified.}
% CAUSALSMITH-CITED-SCOPE-END thm:log-free-bivariate-capacity

\subsection{An independent-prior converse}

The converse in \cref{obj:thm:full-design-lower} certifies the lower comparison scale \leanref{sym:b_s}{\ensuremath{b_s(n,d)}} for every admissible assignment law, including laws that depend on the entire observed covariate sample. The frequency cutoff lets the prior supply the subcritical lower bound as well as the saturated bound. Because the guarantee concerns an average over legal prognostic functions drawn independently of the sample, each admissible design has worst-case loss at least as large as that average.

The next theorem isolates the prior-average lower bound that drives the converse for the entire admissible design class.

\begin{theoremv}{thm:full-design-lower}[Full-design capacity lower bound]*
For every \(s\in(0,1]\), define the lower constant \(c_s\) and the comparison scale \(b_s(n,d)\) by
\[
c_s=\frac{2}{(48+32\pi^2)\,16384^s},
\qquad
b_s(n,d)=\min\left\{1,\left(\frac{\binom d2}{n}\right)^s\right\}.
\]
Then \(c_s>0\), independently of \(n,d\). For all integers \(n,d\geq2\), every design \(\pi\in\mathcal D_{n,d,s}\) of \cref{obj:def:design-class} satisfies
\[
\mathbb E_\xi\,\mathcal L_{n,d,s}(\pi,m_\xi)
\geq c_s b_s(n,d),
\]
where the expectation is over the independent coefficient-sign cosine prior and \(m_\xi\) is its associated prognostic function; the coefficient signs are drawn independently of the covariate sample \(X\). Consequently, the minimax criterion of \cref{obj:def:criterion} satisfies
\[
R_s(n,d)\geq c_s b_s(n,d).
\]
\end{theoremv}
% CAUSALSMITH-CITED-SCOPE-BEGIN thm:full-design-lower
\verificationfootnotetext{\textbf{Formalization scope.} This result uses the cited conclusion \citep{BartlettMendelson2002Complexities} (Theorem 12(4), with Rademacher complexity as in Definition 2). The cited source proof is not formalized here: Lean verifies its use and all remaining steps. If that cited conclusion is false, the portions of this result that depend on it are not certified.}
% CAUSALSMITH-CITED-SCOPE-END thm:full-design-lower

\subsection{Causal transfer and prospective precision}

The completion in \cref{obj:synth_4} retains \(m\) as the conditional midpoint mean while fixing the treatment-effect function and the conditional noise variances. Its scaled variance benchmark is \(4.01\), so a signing-loss bound translates directly into an actual excess-variance bound. This translation uses the excess-variance identity of \citet[Proposition~2.3 (Excess Variance), equation~(2.4), with Assumption~2.1 and Definition~2.2]{Cytrynbaum2026Limits}, under conditional fairness, outcome-independent assignment, and finite potential-outcome second moments.

The exact minimax equality has a causal interpretation. The published excess-variance identity makes the conditional midpoint mean the function that determines excess above the law-specific benchmark; the Gaussian completion realizes every function in the prognostic class. Accordingly, \cref{obj:thm:log-free-bivariate-capacity} equates the prognostic minimax criterion with the outcome-law minimax criterion over \(\mathcal Q_{d,s}\), and gives the corresponding equality for the worst-case loss of the specified procedure. The same capacity is realized within the Gaussian subclass, where the benchmark is \(4.01\).

For prospective precision calculations under the Gaussian completion, the sufficient sample-size bound guarantees scaled variance at most \(4.01(1+\varepsilon)\) uniformly over the prognostic class. The necessary bound applies when \(0<\varepsilon<c_s/4.01\). At \(s=1\), \(c_1\approx3.36\times10^{-7}\), so this displayed tolerance range is approximately \(\varepsilon<8.4\times10^{-8}\); the ratio \(3072/c_1\) exceeds nine billion. These constants are conservative finite-sample certificates of the interaction-count scaling rather than a claim of sharp numerical calibration.

The larger bounded-density class in \cref{obj:lem:published-prognostic-capacity} inherits the converse because its minimax scaled excess is at least \(R_s(n,d)\). For fixed density bounds \(0<\mathrm{lower}\leq1\leq\mathrm{upper}\) and fixed smoothness, vanishing minimax scaled excess in that class therefore requires \(d_n=o(\sqrt n)\). This necessary dimension condition applies with the permitted intercept, the stated prognostic second-moment bound, and finite potential-outcome second moments.
